\section{Demostraciones complementarias de SP7}
\label{app:sp7-proofs}

Estas pruebas suponen catálogos finitos, recursos no ponderados, costes fijos, mensajes coherentes y movimiento muestreado sobre un grafo inflado. No cubren dinámica continua, errores de localización, contacto ni red imperfecta.

\subsection{Identidad de potencial exacto}
\label{app:sp7-potential-proof}

Fijadas las rutas $\boldsymbol r_{-i}$, sea $m_e^{-i}$ el número de ellas que usa el recurso $e$. Entonces
\[
 U_i(r_i,\boldsymbol r_{-i})=-b_{ir_i}^7-\lambda_7\sum_{e\in\mathcal E_{ir_i}^7}m_e^{-i}.
\]
Como $\binom{m+1}{2}-\binom m2=m$, añadir la ruta de $i$ incrementa cada término de pares exactamente en $m_e^{-i}$. Para cualquier cambio $r_i\to r_i'$ se obtiene por tanto
\[
 \Phi_7(r_i',\boldsymbol r_{-i})-\Phi_7(r_i,\boldsymbol r_{-i})
 =U_i(r_i',\boldsymbol r_{-i})-U_i(r_i,\boldsymbol r_{-i}).
\]
El juego es de potencial exacto \parencite{monderer1996potential}. Su espacio tiene $\prod_i|\mathcal R_i^7|$ perfiles. Toda mejora estricta aumenta $\Phi_7$, no repite perfil y termina en un Nash tras, como máximo, $\prod_i|\mathcal R_i^7|-1$ cambios. Esto prueba el Teorema~\ref{thm:sp7-exact-potential}.

\subsection{Condición suficiente para excluir conflictos}
\label{app:sp7-conflict-free-proof}

Si $P_7(\boldsymbol r)>0$ y $\theta_7<\infty$, la definición~\eqref{eq:sp7-threshold} proporciona una desviación con $\Delta P=P_7(\boldsymbol r)-P_7(a,\boldsymbol r_{-i})>0$ y $[\Delta B]_+/\Delta P\leq\theta_7$, donde $\Delta B=B_7(a,\boldsymbol r_{-i})-B_7(\boldsymbol r)$. Su variación de potencial es $-\Delta B+\lambda_7\Delta P>0$ para $\lambda_7>\theta_7$: es inmediata si $\Delta B\leq0$ y, en otro caso, se sigue de $\lambda_7\Delta P>\theta_7\Delta P\geq\Delta B$. Por exactitud, el perfil conflictivo no es Nash. Todo Nash satisface $P_7=0$, como afirma la Proposición~\ref{prop:sp7-conflict-free}. Si no existe desviación unilateral que reduzca $P_7$, entonces $\theta_7=\infty$ y la conclusión no aplica. Esta es la frontera de los pasillos y cuellos de botella.

\subsection{Invariante lógico de la reserva}
\label{app:sp7-reservation-invariant}

Inicialmente cada zona está libre. El arbitraje asigna un único propietario y rechaza entradas ajenas mientras exista su testigo. Este solo se libera después de observar la salida del propietario y cumplir el intervalo de despeje. Por inducción sobre los pasos, nunca se autorizan dos entradas incompatibles. Además, el ejecutor conserva una propuesta por destino, rechaza intercambios opuestos y movimientos hacia agentes estacionarios. Si los nodos activos son distintos al inicio, también lo son al final. La exclusión lógica queda invariante.

\subsection{Ausencia condicional de inanición}
\label{app:sp7-no-starvation}

Sea $S$ un conjunto finito de solicitantes persistentes. Mientras no avanza, $w_i$ aumenta. Tras cruzar, el propietario libera el testigo y abandona la competencia. En cada liberación se elige el máximo lexicográfico de $(w_i,P_i,-i)$. Si algún $j\in S$ nunca fuese elegido, las concesiones retirarían en tiempo finito a los demás hasta dejarlo como máximo, contradicción. El argumento falla si un propietario no libera, aparecen solicitantes sin límite o un agente reingresa indefinidamente.

Estos invariantes excluyen coincidencias de celda, no certifican distancia euclídea entre muestras. Una carga extensa aún puede violarla si la interpolación, la orientación o la dinámica no respetan la inflación. La ejecución física requiere una guardia como la de SP5.
